% !TeX program = tectonic
\documentclass[11pt,a4paper]{article}
\usepackage{fontspec}
\usepackage[english]{babel}
\babelfont{rm}[Language=Default]{Liberation Serif}
\babelfont[Language=Default]{sf}{Carlito}
\babelfont[Language=Default]{tt}{DejaVu Sans Mono}
\usepackage{amsmath,amssymb,amsthm}
\usepackage{geometry}
\geometry{a4paper, top=2.4cm, bottom=2.4cm, left=2.4cm, right=2.4cm}
\usepackage{booktabs,tabularx,enumitem,xcolor,tcolorbox}
\tcbuselibrary{breakable,skins}
\definecolor{linkblue}{HTML}{1F4E79}
\definecolor{statgreen}{HTML}{2F6B3A}
\definecolor{goldacc}{HTML}{8B7E5A}
\usepackage[colorlinks=true,linkcolor=linkblue,urlcolor=linkblue,unicode=true]{hyperref}
\theoremstyle{plain}
\newtheorem{theorem}{Theorem}
\newtheorem{lemma}[theorem]{Lemma}
\newtheorem{proposition}[theorem]{Proposition}
\newtheorem{corollary}[theorem]{Corollary}
\theoremstyle{definition}
\newtheorem{definition}[theorem]{Definition}
\theoremstyle{remark}
\newtheorem{remark}[theorem]{Remark}
\newtcolorbox{statusbox}[1][]{colback=statgreen!5,colframe=statgreen!75,
  title={\textbf{Summary of results:} #1},fonttitle=\small,boxrule=0.7pt,arc=2pt,breakable}
\newtcolorbox{databox}[1][]{colback=goldacc!6,colframe=goldacc!80,
  title={\textbf{Protocol data:} #1},fonttitle=\small,boxrule=0.7pt,arc=2pt,breakable}
\newtcolorbox{verifybox}[1][]{colback=linkblue!5,colframe=linkblue!80,
  title={\textbf{Verification:} #1},fonttitle=\small,boxrule=0.7pt,arc=2pt,breakable}
\widowpenalty=10000
\clubpenalty=10000
\setlength{\parskip}{0.35em}
\newcommand{\CC}{\mathbb{C}}
\newcommand{\RR}{\mathbb{R}}
\newcommand{\ZZ}{\mathbb{Z}}
\newcommand{\QQ}{\mathbb{Q}}
\newcommand{\Ga}{\Gamma}
\newcommand{\Bfun}{\mathrm{B}}
\newcommand{\im}{\operatorname{Im}}
\newcommand{\re}{\operatorname{Re}}
\newcommand{\lcm}{\operatorname{lcm}}
\newcommand{\SNF}{\operatorname{SNF}}
\newcommand{\Jac}{\operatorname{Jac}}
\newcommand{\rank}{\operatorname{rank}}
\newcommand{\Ind}{\operatorname{Index}}
\newcommand{\disc}{\operatorname{disc}}
\begin{document}
\begin{center}
{\LARGE\bfseries The Torus Triple and the Corrections}\\[0.4em]
{\large the calibration stand from first principles}\\[0.6em]
{\normalsize Isaev Iskhak Khamzatovich}\\[0.2em]
{\normalsize The <<Dynamic Principle>> Program --- the \texttt{hodge-laboratory}}\\[0.15em]
{\small Standalone theorem monograph · Монография 7/16}
\end{center}
\vspace{0.4em}
{\color{linkblue}\hrule height 0.8pt}
\vspace{0.8em}

\section{Setting}

The torus is the calibration stand: its triple $(4,1,4\pi^2)$ and
all derived quantities follow from first principles, with no
external premise. The whole correction pipeline is visible at once,
and any construction error would show at coarse precision.

\begin{lemma}[the torus triple]\label{lem:torus}
The torus supplies the integer triple $(4,1,4\pi^2)$: rotation
order $n=4$, carrier index $B=1$, spectral threshold
$\lambda_0=4\pi^2$.
\end{lemma}

\begin{proof}
The rotation $x\mapsto ix$ has order $4$: $i^4=1$, $i^2=-1\ne1$.
The carrier is the torus itself: $H_2(T^2,\ZZ)\cong\ZZ$. The flat
Laplacian spectrum is $|m\tau-n|^2$; in the Tate normalization the
threshold is $4\pi^2$ (certificate G1).
\end{proof}

\begin{theorem}[the corrections derived on the torus]\label{th:toruscorr}
From $(4,1,4\pi^2)$ the pipeline rules give:
$\delta=\pi/4=0.7853981634$; the rotation holonomy $=i$ (protocol
deviation $6.1\cdot10^{-17}$ --- machine precision);
$\delta^2/2=0.3084251375$; $\gamma=\delta^4=0.3805042619$;
$\delta_{\mathrm{eff}}=\delta^5=0.2988473484$;
$\Delta_{Ch}^{\mathrm{torus}}=39.4879953935$.
\end{theorem}

\begin{proof}
Each quantity follows one rule: $\delta=\pi/N$ (the spin double
cover); $\delta^2/2$ --- the main phase correction;
$\gamma=\delta^4/k$ at $k=1$; $\delta_{\mathrm{eff}}
=\delta\cdot\gamma$; the assembly by $\Delta_{Ch}$ at $R=0$. The
holonomy: the lift of the rotation permutes the sheets; the
eigenvalue on an eigenform is $i$. All digits reproduce at $dps=35$
with zero last-digit discrepancy.
\end{proof}

\begin{remark}[chessboard equivariance]
For the symmetric encoding the $\mu_4$-equivariance is exact
(deviation $0.000$ over the $4096\times4096$ grid); for the $(7,22)$
encoding it is broken ($1.000$). The contrast proves the corrections
are geometric --- they inherit the carrier symmetry.
\end{remark}

\begin{databox}[numbers]
The phase family $\delta_T=\pi/N$ reproduced for $N=2,\dots,8$; the
$\gamma$, $\delta_{\mathrm{eff}}$ tables for $k\in\{1,22\}$ built in
full; the E6 case corrected: $\Delta=0.7990$ (not $0.5964$); on
Klein the same formula gives $\Delta_{Ch}\approx3.4399$ against the
observed $3.443$ --- deviation $0.0905\%$.
\end{databox}

\begin{statusbox}[summary]
The torus triple and all derivatives are derived from first
principles; the exact torus run calibrates all higher rungs.
\end{statusbox}

\begin{verifybox}[reproduction]
\texttt{python3 laboratory.py} $\to$ ``Stands $\to$ Torus''; the
plot \texttt{fig\_torus.png}.
\end{verifybox}

\vfill
{\color{linkblue}\hrule height 0.6pt}
\vspace{0.5em}
{\small\color{gray}
License: individual exclusive license of Isaev Iskhak Khamzatovich.
All rights to the computations, formulas, and texts belong to the author.
Verification: \texttt{python3 laboratory.py} (``Stands''/``Certificates''),
Lean: \texttt{verification/lean/HodgeLaboratory.lean}.
}
\end{document}
